\documentclass[11pt]{article}
\usepackage{fontspec}
\usepackage{xeCJK}
\setCJKmainfont[ItalicFont=STSongti-SC-Bold]{Songti SC}
\usepackage{xurl}
\usepackage[a4paper,margin=25mm]{geometry}
\usepackage{amsmath,amssymb,amsthm,mathtools}
\usepackage{xcolor}
\usepackage{booktabs,array,enumitem,fancyhdr,hyperref}
\hypersetup{colorlinks=true,linkcolor=blue!50!black,urlcolor=blue!50!black}
\setlength{\emergencystretch}{2em}
\setlength{\parskip}{0.3em}
\linespread{1.16}
\allowdisplaybreaks
\newtheorem{theorem}{定理}[section]
\newtheorem{proposition}[theorem]{命题}
\newtheorem{lemma}[theorem]{引理}
\theoremstyle{definition}
\newtheorem{definition}[theorem]{定义}
\theoremstyle{remark}
\newtheorem{remark}[theorem]{说明}
\newcommand{\dd}{\,\mathrm d}
\newcommand{\ii}{\mathrm i}
\newcommand{\Rea}{\operatorname{Re}}
\newcommand{\Ima}{\operatorname{Im}}
\newcommand{\Li}{\operatorname{Li}}
\newcommand{\supp}{\operatorname{supp}}
\pagestyle{fancy}
\fancyhf{}
\fancyhead[L]{\small 学习笔记}
\fancyhead[R]{\small 初稿}
\fancyfoot[C]{\thepage}
\setlength{\headheight}{15pt}
\author{}
\date{2026年10月22日初稿}
\makeatletter
\renewcommand{\maketitle}{%
  \thispagestyle{plain}%
  \begin{center}
    {\LARGE\bfseries \@title\par}
    \vspace{0.6em}
    {\small \@date\par}
  \end{center}
  \vspace{0.6em}}
\makeatother
\newcommand{\draftnotice}{\noindent{\small
本文为供逐篇审阅的学习笔记初稿，不是新增研究成果。
文中区分直接证明与引用结果；正文不依赖本文。
本轮已作公式、条件及排版初步核查，尚待作者逐句审定。}\par\medskip}

\XeTeXlinebreaklocale "zh"
\renewcommand{\proofname}{证明}
\DeclareMathOperator{\Res}{Res}
\DeclareMathOperator{\PV}{p.v.}
\date{2026年10月24日初稿}
\title{极点为何在复变函数中如此有用}
\begin{document}
\maketitle

\noindent{\small 在实变函数中，使函数无定义的点通常从定义域中剔除，只作为需要回避的例外；
在复变函数中，同样的点却成为计算积分、求级数和、确定渐近行为的主要工具。
本篇说明这一转变的数学原因。凡标为命题者给出证明；标准定理注明出处后引用。}\par\medskip

\section{问题的确切含义}
需要先说明：极点在复平面上同样没有函数值，并没有因为进入复数域而“被定义”。
发生变化的是极点周围的结构。本篇论证以下四点：
\begin{enumerate}
 \item 全纯性是很强的约束，使函数在极点附近只能有有限多个自由参数（第 2 节）；
 \item 平面上存在绕孤立点的闭曲线，积分因而能够探测奇点（第 3 节）；
 \item 实解析函数是复解析函数的限制，复平面上“看不见”的极点控制实轴上的行为（第 4 节）；
 \item 在适当的整体条件下，函数由其极点的数据决定（第 5 节）。
\end{enumerate}
第 6 节说明复视角如何为实轴上的奇异积分赋予意义，第 7 节说明这些原因在本书显式公式中的体现。

\section{刚性：极点附近只有有限个参数}
\subsection{实变量下奇点附近没有约束}
设 $f$ 在 $(-1,0)\cup(0,1)$ 上光滑。两个区间互不连通，$f$ 在两侧的取值彼此独立：
例如可以令 $f(x)=1/x$（$x>0$），$f(x)=0$（$x<0$），也可以把右侧换成任何别的光滑函数。
即使 $f$ 在 $0$ 附近有界，也未必能连续延拓到 $0$：$\sin(1/x)$ 就是例子。
因此实变量下，“$0$ 是奇点”这一事实本身几乎不包含信息。

\subsection{复变量下的 Laurent 展开}
复平面中，去心圆盘 $0<|z-a|<r$ 是连通的。若 $f$ 在其中全纯，则有 Laurent 展开
\[
 f(z)=\sum_{n=-\infty}^{\infty}c_n(z-a)^n,\qquad
 c_n=\frac1{2\pi\ii}\oint_{|z-a|=\rho}\frac{f(z)}{(z-a)^{n+1}}\dd z\quad(0<\rho<r),
\]
展开唯一，且系数与 $\rho$ 无关（Ahlfors, \emph{Complex Analysis}, Ch.~5）。孤立奇点只有三类：
\begin{itemize}
 \item 若 $f$ 在 $a$ 附近有界，则 $a$ 是可去奇点（Riemann 可去奇点定理）。
 \item 若只有有限多个负幂项，$a$ 是\emph{极点}，$f$ 在 $a$ 附近由有限的\emph{主部}
 $\sum_{n=1}^mc_{-n}(z-a)^{-n}$ 加一个全纯函数构成。
 \item 否则 $a$ 是本性奇点，$f$ 在 $a$ 的任意邻域内的值在 $\mathbb C$ 中稠密（Casorati--Weierstrass）。
\end{itemize}
与实变量对照：$\sin(1/x)$ 在实轴上有界，而 $\sin(1/z)$ 在 $0$ 附近无界，$0$ 是本性奇点。
有界性在复变量下是强得多的条件。

于是，极点附近的全部“奇异信息”集中在主部的有限多个系数中；
其余部分是全纯的，由恒等定理在整个连通区域上相互牵连。
这种刚性是极点能够携带信息的前提：一个可以任意修改的点不携带信息，一个由有限参数刻画、且与整体相联的点才携带信息。

\section{积分：留数是拓扑信息}
\begin{proposition}\label{prop:residue}
设 $f$ 在 $0<|z-a|<r$ 内全纯，Laurent 系数为 $c_n$。则对 $0<\rho<r$，
\[
 \oint_{|z-a|=\rho}f(z)\dd z=2\pi\ii\,c_{-1}.
\]
\end{proposition}
\begin{proof}
Laurent 级数在圆周上一致收敛，可以逐项积分。对整数 $n$，
$\oint_{|z-a|=\rho}(z-a)^n\dd z=\int_0^{2\pi}\rho^{n+1}e^{\ii(n+1)\theta}\ii\dd\theta$，
当 $n\ne-1$ 时为零，$n=-1$ 时为 $2\pi\ii$。
\end{proof}
记 $\Res_{z=a}f=c_{-1}$。由 Cauchy 定理，对区域内的闭曲线，积分只依赖于曲线所围的奇点及各自的留数（留数定理）。

这一事实的本质是拓扑的：在去心平面上，$\dd z/z$ 是闭形式而不是恰当形式，
它沿闭曲线的积分等于 $2\pi\ii$ 乘以绕数。实直线去掉一点后成为两个区间，不存在绕该点的闭曲线；
实变量下没有任何积分能够“包围”一个奇点。
因此，只有在二维的复平面中，奇点才能通过积分被测量。

\begin{proposition}\label{prop:fourier}
对实数 $\xi$，$\displaystyle\int_{-\infty}^{\infty}\frac{e^{-\ii x\xi}}{1+x^2}\dd x=\pi e^{-|\xi|}$。
\end{proposition}
\begin{proof}
设 $\xi\ge0$。在下半平面取半圆周 $|z|=R$，$\Ima z\le0$，$R>1$。在其上 $|e^{-\ii z\xi}|=e^{\xi\Ima z}\le1$，
$|1+z^2|\ge R^2-1$，故半圆上的积分不超过 $\pi R/(R^2-1)\to0$。
闭曲线按顺时针方向包围唯一的极点 $z=-\ii$，留数为 $e^{-\ii(-\ii)\xi}/(-2\ii)=e^{-\xi}/(-2\ii)$，
故积分为 $-2\pi\ii\cdot e^{-\xi}/(-2\ii)=\pi e^{-\xi}$。$\xi<0$ 时在上半平面同理处理。
\end{proof}
这个结果的形式也值得注意：Fourier 变换的指数衰减率 $1$，恰是极点 $\pm\ii$ 到实轴的距离。
一般地，函数在实轴附近带形区域中的解析性决定其 Fourier 变换的衰减速度（Paley--Wiener 型定理）。

\section{看不见的极点控制实轴上的行为}
\subsection{收敛半径}
$1/(1+x^2)$ 在整条实轴上无穷次可微，但它在 $0$ 处的 Taylor 级数 $\sum(-1)^nx^{2n}$ 的收敛半径只有 $1$。
仅从实轴无法解释这个“$1$”。复平面上的解释是：幂级数的收敛半径等于展开中心到最近奇点的距离
（Cauchy--Hadamard 公式与解析延拓；Ahlfors, Ch.~5），而 $1/(1+z^2)$ 在 $\pm\ii$ 处有极点，到 $0$ 的距离为 $1$。

\subsection{系数的渐近：生成函数}
\begin{proposition}\label{prop:fib}
设 $F_0=0$，$F_1=1$，$F_{n+1}=F_n+F_{n-1}$。记 $\varphi=\frac{1+\sqrt5}2$，$\bar\varphi=-1/\varphi$。则
$F_n=(\varphi^n-\bar\varphi^n)/\sqrt5$。
\end{proposition}
\begin{proof}
由递推，$F(z)=\sum_{n\ge0}F_nz^n=z/(1-z-z^2)$，在 $|z|<1/\varphi$ 内收敛。
$1-z-z^2=-(z-a)(z-b)$，$a=1/\varphi$，$b=-\varphi$。对 $n\ge0$，
$F(z)z^{-n-1}$ 的留数之和由三项组成：$z=0$ 处为 $F_n$；$z=a$ 处为 $-a^{-n}/(a-b)$；$z=b$ 处为 $-b^{-n}/(b-a)$。
在 $|z|=R$ 上，$F(z)z^{-n-1}=O(R^{-n-2})$，所以大圆周上的积分在 $R\to\infty$ 时趋于零，留数之和为零。
由 $a^{-n}=\varphi^n$，$b^{-n}=\bar\varphi^n$，$a-b=\sqrt5$ 即得。
\end{proof}
证明的结构具有一般性：系数由生成函数在各极点的留数给出，距原点最近的极点给出主导项，较远的极点给出较小的修正。
解析组合学中的“奇点分析”把这一原理推广到代数型与对数型奇点（Flajolet--Sedgewick, \emph{Analytic Combinatorics}, Ch.~IV, VI）。

\section{函数由极点决定}
在整体条件下，极点的数据可以完全确定函数。
\begin{proposition}\label{prop:sin}
对 $z\notin\mathbb Z$，$\displaystyle\frac{\pi^2}{\sin^2\pi z}=\sum_{n\in\mathbb Z}\frac1{(z-n)^2}$。
特别地，$\sum_{n\ge1}n^{-2}=\pi^2/6$。
\end{proposition}
\begin{proof}
右端在不含整数的紧集上一致收敛，定义一个以 $1$ 为周期的亚纯函数，在每个整数 $n$ 处有主部 $(z-n)^{-2}$。
由 $\sin\pi z=\pi z\bigl(1-\tfrac{\pi^2z^2}6+O(z^4)\bigr)$，
\[
 \frac{\pi^2}{\sin^2\pi z}=\frac1{z^2}+\frac{\pi^2}3+O(z^2),                  \tag{5.1}
\]
左端在 $0$ 处的主部也是 $z^{-2}$；由周期性，在每个整数处亦然。
于是差 $g(z)$ 是以 $1$ 为周期的整函数。在带形 $0\le\Rea z\le1$ 中，当 $|\Ima z|\to\infty$ 时，
$|\sin\pi z|\ge\sinh(\pi|\Ima z|)$ 使左端趋于零，右端每项不超过 $1/((\Rea z-n)^2+(\Ima z)^2)$，其和也一致趋于零。
因此 $g$ 有界；由 Liouville 定理 $g$ 为常数，而该常数是上述极限 $0$。
最后，在恒等式中减去 $n=0$ 项并令 $z\to0$，由 (5.1) 得 $2\sum_{n\ge1}n^{-2}=\pi^2/3$。
\end{proof}
证明中用到的只有两件事：两边的主部相同，以及在无穷远处的增长受控。
一般地，Riemann 球面上的亚纯函数都是有理函数，由其极点的主部确定到差一个常数；
在平面上，Mittag-Leffler 定理允许任意指定极点及主部，Weierstrass 乘积定理允许任意指定零点（Ahlfors, Ch.~5）。
在这一意义上，亚纯函数“由其奇点构成”。

\section{复视角为实轴上的奇异积分赋予意义}
$\int\varphi(x)/x\dd x$ 在 $0$ 处一般发散。复视角给出一种有确定意义的处理：把奇点从实轴移开。
\begin{proposition}[Sokhotski--Plemelj]\label{prop:plemelj}
设 $\varphi$ 在 $\mathbb R$ 上连续可微且有紧支集。则
\[
 \lim_{\varepsilon\to0^+}\int_{-\infty}^{\infty}\frac{\varphi(x)}{x-\ii\varepsilon}\dd x
 =\PV\!\int_{-\infty}^{\infty}\frac{\varphi(x)}x\dd x+\ii\pi\varphi(0),
\]
其中 $\PV\!\int\varphi(x)/x\dd x=\lim_{\delta\to0^+}\int_{|x|>\delta}\varphi(x)/x\dd x=\int_0^\infty\frac{\varphi(x)-\varphi(-x)}x\dd x$。
\end{proposition}
\begin{proof}
$\dfrac1{x-\ii\varepsilon}=\dfrac{x}{x^2+\varepsilon^2}+\ii\dfrac{\varepsilon}{x^2+\varepsilon^2}$。
实部：$\int\varphi(x)\frac{x}{x^2+\varepsilon^2}\dd x=\int_0^\infty(\varphi(x)-\varphi(-x))\frac{x}{x^2+\varepsilon^2}\dd x$；
由 $|\varphi(x)-\varphi(-x)|\le2x\sup|\varphi'|$ 及紧支性，控制收敛定理给出主值积分。
虚部：令 $x=\varepsilon t$，$\int\varphi(x)\frac{\varepsilon}{x^2+\varepsilon^2}\dd x=\int\frac{\varphi(\varepsilon t)}{1+t^2}\dd t\to\pi\varphi(0)$。
\end{proof}
实轴上的奇点在这里分解为两部分：主值部分对应两侧的对称相消，
$\ii\pi\varphi(0)$ 对应半个留数——从下方绕过极点的小半圆贡献 $\pi\ii$ 而不是 $2\pi\ii$。
实变量下只能“回避”的奇点，在复平面上得到确定的贡献。

\section{在本书中的体现}
本书的主要公式几乎都是留数的和：
\begin{itemize}
 \item Perron 公式把 $\psi(x)$ 写成 $\frac1{2\pi\ii}\int_{(c)}\bigl(-\frac{\zeta'}{\zeta}(s)\bigr)\frac{x^s}s\dd s$。
 左移积分线后，每个极点贡献一项：
 \begin{center}
 \small
 \begin{tabular}{lll}
 \toprule
 极点 & 来源 & 显式公式中的项\\
 \midrule
 $s=1$ & $\zeta$ 的极点 & 主项 $x$\\
 $s=\rho$ & $\zeta$ 的非平凡零点 & $-x^\rho/\rho$\\
 $s=0$ & 核 $1/s$ & $-\log2\pi$\\
 $s=-2n$ & $\zeta$ 的平凡零点 & $-x^{-2n}/(-2n)$\\
 \bottomrule
 \end{tabular}
 \end{center}
 \item 零点进入公式的方式本身就是极点：$\zeta$ 的零点是对数导数 $\zeta'/\zeta$ 的单极点，留数等于零点重数。
 取对数导数这一步，把“函数在哪里为零”转化为“哪里有极点、留数多少”，从而可以用留数定理计算。
 \item 学习笔记第 11 篇中 $\zeta(-1)=-\tfrac1{12}$ 与 $\zeta$ 在 $s=1$ 处的极点，学习笔记第 8 篇中 $\psi_1$ 公式多出的 $s=-1$ 极点，
 都是同一原理的例子：核的选择决定出现哪些极点，极点决定公式中出现哪些项。
\end{itemize}

\section{讨论}
综合以上各节，极点在复变函数中有用，原因不在于极点本身，而在于全纯性施加的约束：
\begin{enumerate}
 \item 全纯函数是刚性的：局部决定整体，奇点附近只允许有限参数（主部）。
 \item 复平面是二维的：存在绕孤立点的闭曲线，积分可以测量奇点（留数），这是一个拓扑不变量。
 \item 实解析函数是复解析函数的限制：实轴上的收敛半径、Fourier 衰减与系数渐近，由复平面中的极点决定。
 \item 在整体增长受控时，函数由奇点的数据决定（Liouville、Mittag-Leffler），因而“知道极点就知道函数”。
\end{enumerate}
在约束很强的理论中，约束失效之处携带了全部信息。
\emph{旁注}：静电学中的调和函数与此类似，在已知边界条件下由其源（点电荷）决定，
二维中 $\log|z-a|$ 的电荷正对应 $1/(z-a)$ 的留数；这一类比只说明结构上的相似。
实变函数中奇点之所以“无用”，是因为实变量没有提供这种约束，奇点周围的函数可以任意改变。
\end{document}
